\par\noindent\textbf{(1)} 当地竖直的地球自转分量为 Ω\_z=Ω\_E sinφ，重力提供回复频率 ω\_0=√(g/l)。科氏加速度为 −2Ω×v，投影得 ẍ+ω\_0²x=2Ω\_z ẏ、ÿ+ω\_0²y=−2Ω\_z ẋ，符号对应 x 向东 y 向北的右手约定。
\par\noindent\textbf{(2)} 取 z=x+iy，则 z̈+2iΩ\_z ż+ω\_0²z=0。代入 z∝e\textasciicircum{}\{−iωt\} 得 −ω²+2Ω\_zω+ω\_0²=0，两个带符号根为 Ω\_z±√(ω\_0²+Ω\_z²)。实际两圆偏振频率模大小分裂约 2Ω\_z。
\par\noindent\textbf{(3)} 把解写成 z=e\textasciicircum{}\{−iΩ\_z t\}[A cos(√(ω\_0²+Ω\_z²)t)+B sin(√(ω\_0²+Ω\_z²)t)]。快速摆动的线偏振方向随因子 e\textasciicircum{}\{−iΩ\_z t\} 进动，故相对地面角速度为 −Ω\_E sinφ，周期 T\_p=2π/(Ω\_E|sinφ|)，赤道处周期趋于无穷。
\par\noindent\textbf{(4)} 极点 sinφ=±1，进动一恒星日一周；赤道 sinφ=0，不出现主阶岁差。要求 Ω\_E≪√(g/l)，摆角≪1，且地球本地坐标和悬点近似静止。真正地球摆还受非理想椭圆轨迹、支承扭矩和空气阻尼影响，不能把此简化结论误认精确实验周期。
